\chapter{Discussion, limitations, and conclusion}\label{sec:conclusion}
A simple deterministic local toppling rule produces broad avalanche responses as inputs encounter different configurations. The baseline has unequal stable-height frequencies and finite-grid mean topplings that agree with analytic predictions. Size, area and reach have closely related ranks, while loss depends more on access to a sink.

The three extensions distinguish changes in state, response and event definition. On a fixed graph, full-support driving preserves the stationary recurrent measure but changes how input sites weight the response. Cylindrical boundaries alter connectivity and dissipation, changing the measured stationary heights and mean topplings. Their response is governed by the inverse toppling matrix $K^{-1}$, which maps mean input to mean sitewise topplings in Eq.~\eqref{eq:benchmark}. Batching changes event distributions while preserving grouped sequential toppling totals, loss totals and final configurations exactly. Its larger events combine successive, potentially overlapping responses.

The evidence is limited to four baseline sizes, two sizes per extension and five primary seeds per condition. Temporal dependence motivates seed-level and block-based uncertainty estimates, whose precision is limited by the replicate count. Stationarity diagnostics concern selected observables; the confirmation examines one height discrepancy. Finite-range marginal fits have systematic shape errors and range-dependent parameters, so their exponents describe these samples. KS and score comparisons are descriptive, without calibrated IID goodness-of-fit $p$-values. Batches relax completely between inputs, which leaves finite-rate driving untested. Stabilization benchmarks describe one workload; the production runtime measures the complete run.

Conservation fixes long-run mean loss at one per input grain, despite changes in loss frequencies, magnitudes and event grouping. Growing finite-size ranges and better cutoff descriptions support broad, size-limited responses; range-dependent exponents and residual shape do not establish universal power laws. A changed avalanche distribution must therefore be interpreted through the stationary configuration measure, the response to driving and the definition of an event.
